% Part: many-valued-logic
% Chapter: three-valued-logics
% Section: introduction

\documentclass[../../../include/open-logic-section]{subfiles}

\begin{document}

\olfileid{mvl}{thr}{int}

\olsection{परिचय}

$\True$ आणि $\False$ या मूल्यांच्या जोडीला आणखी एक मूल्य~$\Undef$
जोडले, तर त्रिमूल्य तर्कशास्त्र मिळते. हे एकच अतिरिक्त
सत्यतामूल्य असले, तरी $\lnot$, $\land$, $\lor$ आणि $\lif$
यांसाठी सत्यता-फलने परिभाषित करण्याचे अनेक पर्याय असतात.
मग एखादे तर्कशास्त्र या सत्यता-फलनांचा कोणताही संयोग वापरू
शकते; तसेच फक्त $\True$ ला निर्दिष्ट मानायचे की $\True$
आणि~$\Undef$ या दोन्हींना, हाही पर्याय असतो.

येथे आपण सर्वाधिक परिचित त्रिमूल्य तर्कशास्त्रांपैकी काही,
त्यांमागील प्रेरणा आणि त्यांचे काही गुणधर्म मांडतो.

\end{document}
